\chapter{Arrow's Impossibility Theorem}
%%%%%%%%%%%%%%%%%%%%%%%%%%%%%%%%%%%%%%%%%%%%%%%%%%%%%%%%%%%%%%%%%%%%%%%%%%%%%%%
\section{Overview}

This chapter proves \textbf{Arrow's Impossibility Theorem} via Kalai's
Fourier-analytic approach.  For 3 alternatives $\{a,b,c\}$ and $n$ voters,
a social welfare function is modelled by an odd Boolean function
$f:\{0,1\}^n\to\bbr$ applied to each pairwise comparison.  The main result is:

\begin{quote}
Any odd, $\pm 1$-valued, unanimous, acyclic social welfare function must be a
\emph{dictator}.
\end{quote}

The proof proceeds in three steps:
\begin{enumerate}
  \item Acyclicity forces the Fourier correlation function $\mathrm{corr}(f) = -1/3$.
  \item Since $\mathrm{corr}(f)\ge -1/3$ always holds (with equality iff all
        Fourier weight is on level 1), this forces $W_1[f] = 1$.
  \item A $\pm 1$-valued degree-1 unanimous function must be a dictator.
\end{enumerate}

\textbf{References:}
\begin{itemize}
  \item Gil Kalai, \emph{A Fourier-Theoretic Perspective on the Condorcet
        Paradox and Arrow's Theorem}, Advances in Applied Mathematics, 2002.
  \item Ryan O'Donnell, \emph{Analysis of Boolean Functions}, Chapter 2.
\end{itemize}

%%%%%%%%%%%%%%%%%%%%%%%%%%%%%%%%%%%%%%%%%%%%%%%%%%%%%%%%%%%%%%%%%%%%%%%%%%%%%%%
\section{Social choice definitions}

\begin{definition}[Unanimity]
\lean{ArrowTheorem.unanimity}
\leanok
\uses{BooleanAnalysis.BooleanFunc}
A social welfare function $f:\{0,1\}^n\to\bbr$ is \emph{unanimous} if
$f(\texttt{false},\dots,\texttt{false}) = 1$: when all voters prefer
alternative $a$ over $b$, so does society.
\end{definition}

\begin{definition}[Dictatorship]
\lean{ArrowTheorem.isDictator}
\leanok
\uses{BooleanAnalysis.dictator}
$f$ is a \emph{dictatorship} if there exists a voter $i$ such that
$f = \mathrm{dict}_i$, i.e.\ society's preference always equals voter $i$'s.
\end{definition}

%%%%%%%%%%%%%%%%%%%%%%%%%%%%%%%%%%%%%%%%%%%%%%%%%%%%%%%%%%%%%%%%%%%%%%%%%%%%%%%
\section{Voter ordering model}

Each voter holds one of the 6 strict transitive orderings of $\{a,b,c\}$.
For each ordering we record three pairwise preferences using the sign
convention $\sigma(\texttt{false}) = 1$ (``pro first alternative'') and
$\sigma(\texttt{true}) = -1$ (``pro second'').

The six orderings and their signs $(s_{ab}, s_{bc}, s_{ca})$ are:
\begin{center}
\begin{tabular}{clccc}
Index & Ordering & $s_{ab}$ & $s_{bc}$ & $s_{ca}$ \\\hline
0 & $a>b>c$ & $1$ & $1$ & $-1$ \\
1 & $a>c>b$ & $1$ & $-1$ & $-1$ \\
2 & $b>a>c$ & $-1$ & $1$ & $-1$ \\
3 & $b>c>a$ & $-1$ & $1$ & $1$ \\
4 & $c>a>b$ & $1$ & $-1$ & $1$ \\
5 & $c>b>a$ & $-1$ & $-1$ & $1$ \\
\end{tabular}
\end{center}

\begin{definition}[Pairwise preference functions]
\lean{ArrowTheorem.abPref}
\lean{ArrowTheorem.bcPref}
\lean{ArrowTheorem.caPref}
\leanok
\uses{BooleanAnalysis.boolToSign}
$\mathrm{abPref}(k)$, $\mathrm{bcPref}(k)$, $\mathrm{caPref}(k)$ give the
Boolean preference of ordering $k\in\{0,\dots,5\}$ in the $a$-vs-$b$,
$b$-vs-$c$, and $c$-vs-$a$ comparisons, respectively.
\end{definition}

\begin{definition}[Profile]
\lean{ArrowTheorem.Profile}
\leanok
A \emph{profile} assigns each of the $n$ voters one of the 6 orderings:
$p : \mathrm{Fin}\,n \to \mathrm{Fin}\,6$.
\end{definition}

\begin{definition}[Vote vectors]
\lean{ArrowTheorem.abVotes}
\lean{ArrowTheorem.bcVotes}
\lean{ArrowTheorem.caVotes}
\leanok
\uses{ArrowTheorem.Profile, ArrowTheorem.abPref,
      ArrowTheorem.bcPref, ArrowTheorem.caPref}
Given a profile $p$, the vote vectors $\mathrm{abVotes}(p)$,
$\mathrm{bcVotes}(p)$, $\mathrm{caVotes}(p) \in \{0,1\}^n$ record each
voter's pairwise preference.
\end{definition}

\begin{definition}[Acyclicity]
\lean{ArrowTheorem.acyclic}
\leanok
\uses{ArrowTheorem.abVotes, ArrowTheorem.bcVotes, ArrowTheorem.caVotes}
A social welfare function $f$ is \emph{acyclic} if no profile of transitive
voter orderings produces a Condorcet cycle in society's preferences.
A cycle occurs when
$f(\mathrm{abVotes}(p)) = f(\mathrm{bcVotes}(p)) = f(\mathrm{caVotes}(p)) = 1$
(the cycle $a>b>c>a$) or all equal $-1$ (the reverse cycle).
\end{definition}

%%%%%%%%%%%%%%%%%%%%%%%%%%%%%%%%%%%%%%%%%%%%%%%%%%%%%%%%%%%%%%%%%%%%%%%%%%%%%%%
\section{Key correlation lemmas}

The following lemma encodes the core computation: for a single voter drawn
uniformly from the 6 orderings, the expected product of any two distinct
pairwise preference signs is $-1/3$.

\begin{lemma}[Pairwise sign sums]
\lean{ArrowTheorem.sum_abPref_bcPref}
\lean{ArrowTheorem.sum_bcPref_caPref}
\lean{ArrowTheorem.sum_abPref_caPref}
\leanok
\uses{ArrowTheorem.abPref, ArrowTheorem.bcPref, ArrowTheorem.caPref,
      BooleanAnalysis.boolToSign}
Summing pairwise sign products over all 6 orderings:
\[
  \sum_{k=0}^{5} s_{ab}(k)\cdot s_{bc}(k)
  \;=\; \sum_{k=0}^{5} s_{bc}(k)\cdot s_{ca}(k)
  \;=\; \sum_{k=0}^{5} s_{ab}(k)\cdot s_{ca}(k)
  \;=\; -2.
\]
Hence $\E[s_{ab}\cdot s_{bc}] = -2/6 = -1/3$.
\end{lemma}

%%%%%%%%%%%%%%%%%%%%%%%%%%%%%%%%%%%%%%%%%%%%%%%%%%%%%%%%%%%%%%%%%%%%%%%%%%%%%%%
\section{Fourier formula for the cycle correlation}

\begin{definition}[Correlation function]
\lean{ArrowTheorem.corrFunc}
\leanok
\uses{BooleanAnalysis.fourierCoeff}
The \emph{Fourier correlation function} of $f$ is
\[
  \mathrm{corr}(f) \;=\; \sum_{S\subseteq[n]} \hat f(S)^2\,(-1/3)^{|S|}.
\]
For odd $f$, only odd-level terms contribute.
\end{definition}

\begin{lemma}[Expected pairwise product equals correlation function]
\lean{ArrowTheorem.expected_product_eq_corrFunc}
\leanok
\uses{ArrowTheorem.corrFunc, ArrowTheorem.abVotes, ArrowTheorem.bcVotes,
      BooleanAnalysis.fourierCoeff, BooleanAnalysis.chiS,
      ArrowTheorem.sum_abPref_bcPref}
For voters with i.i.d.\ uniform orderings,
\[
  \E_p\bigl[f(\mathrm{abVotes}(p))\cdot f(\mathrm{bcVotes}(p))\bigr]
  \;=\; \mathrm{corr}(f).
\]
The same identity holds for the $bc$--$ca$ and $ab$--$ca$ pairs.
\end{lemma}

\begin{lemma}[Correlation function lower bound]
\lean{ArrowTheorem.corrFunc_ge_neg_third}
\leanok
\uses{ArrowTheorem.corrFunc, BooleanAnalysis.isOddFunc, BooleanAnalysis.isPmOne,
      BooleanAnalysis.fourierCoeff_odd_even, BooleanAnalysis.parseval_pm_one}
For any odd $\pm 1$-valued function,
\[
  \mathrm{corr}(f) \;\ge\; -1/3.
\]
\emph{Proof sketch:}  For odd $|S|\ge 1$ we have $(-1/3)^{|S|}\ge -1/3$.
Even-level coefficients vanish by oddness.  Hence
$\mathrm{corr}(f) \ge (-1/3)\sum_S \hat f(S)^2 = -1/3$.
\end{lemma}

\begin{lemma}[Equality forces all weight on level 1]
\lean{ArrowTheorem.corrFunc_eq_neg_third_of_weight_one}
\leanok
\uses{ArrowTheorem.corrFunc, ArrowTheorem.corrFunc_ge_neg_third,
      BooleanAnalysis.isOddFunc, BooleanAnalysis.isPmOne,
      BooleanAnalysis.fourierCoeff_odd_even, BooleanAnalysis.parseval_pm_one}
If $\mathrm{corr}(f) = -1/3$, then $\hat f(S) = 0$ for all $S$ with
$|S| \ne 1$.  Combined with oddness ($\hat f(\varnothing) = 0$), all Fourier
weight lies on level 1: $W_1[f] = 1$.
\end{lemma}

%%%%%%%%%%%%%%%%%%%%%%%%%%%%%%%%%%%%%%%%%%%%%%%%%%%%%%%%%%%%%%%%%%%%%%%%%%%%%%%
\section{Acyclicity implies degree-1 structure}

\begin{lemma}[Acyclicity implies $\mathrm{corr}(f) = -1/3$]
\lean{ArrowTheorem.acyclic_implies_corrFunc}
\leanok
\uses{ArrowTheorem.acyclic, ArrowTheorem.corrFunc,
      ArrowTheorem.expected_product_eq_corrFunc,
      BooleanAnalysis.isPmOne}
If $f$ is $\pm 1$-valued and acyclic, then $\mathrm{corr}(f) = -1/3$.

\emph{Proof sketch:}
\begin{enumerate}
  \item For any $\pm 1$ triple $(a,b,c)$ avoiding the two all-equal cycles,
        $ab + bc + ac = -1$.
  \item Summing over all $6^n$ profiles:
        $\sum_p (f_{\!ab}\cdot f_{\!bc} + f_{\!bc}\cdot f_{\!ca} + f_{\!ab}\cdot f_{\!ca}) = -6^n$.
  \item Each pairwise expected product equals $\mathrm{corr}(f)$, so
        $3\cdot\mathrm{corr}(f) = -1$, giving $\mathrm{corr}(f) = -1/3$.
\end{enumerate}
\end{lemma}

%%%%%%%%%%%%%%%%%%%%%%%%%%%%%%%%%%%%%%%%%%%%%%%%%%%%%%%%%%%%%%%%%%%%%%%%%%%%%%%
\section{Degree-1 implies dictatorship}

\begin{lemma}[Degree-1 unanimous $\pm 1$ function is a dictator]
\lean{ArrowTheorem.degree_one_implies_dictator}
\leanok
\uses{ArrowTheorem.isDictator, ArrowTheorem.unanimity,
      BooleanAnalysis.isPmOne, BooleanAnalysis.fourierCoeff,
      BooleanAnalysis.parseval_pm_one, BooleanAnalysis.walsh_expansion,
      BooleanAnalysis.dictator_eq_chi, BooleanAnalysis.chiS_singleton}
Suppose $f$ is $\pm 1$-valued, unanimous, and $\hat f(S) = 0$ for all $|S|\ne 1$.
Then $f$ is a dictator.

\emph{Proof sketch:}  Write $f = \sum_i a_i\chi_{\{i\}}$ where $a_i = \hat f(\{i\})$.
\begin{enumerate}
  \item Parseval: $\sum_i a_i^2 = 1$.
  \item Unanimity: $f(\mathbf{0}) = \sum_i a_i = 1$.
  \item For each $j$, the value $f(e_j) = 1 - 2a_j \in\{-1,1\}$ forces $a_j\in\{0,1\}$.
  \item From $a_j\in\{0,1\}$ and $\sum a_j^2 = \sum a_j = 1$: exactly one
        $a_{j_0} = 1$ and the rest are $0$.
  \item Hence $f = \chi_{\{j_0\}} = \mathrm{dict}_{j_0}$.
\end{enumerate}
\end{lemma}

%%%%%%%%%%%%%%%%%%%%%%%%%%%%%%%%%%%%%%%%%%%%%%%%%%%%%%%%%%%%%%%%%%%%%%%%%%%%%%%
\section{Arrow's Impossibility Theorem}

\begin{theorem}[Arrow's Impossibility Theorem]
\lean{ArrowTheorem.arrow_theorem}
\leanok
\uses{ArrowTheorem.acyclic_implies_corrFunc,
      ArrowTheorem.corrFunc_eq_neg_third_of_weight_one,
      ArrowTheorem.degree_one_implies_dictator}
Let $f:\{0,1\}^n\to\bbr$ be a social welfare function that is odd
(antisymmetric), $\pm 1$-valued, unanimous, and acyclic.  Then $f$ is a
\emph{dictatorship}: there exists a voter $i_0$ such that
$f = \mathrm{dict}_{i_0}$.

\emph{Proof:}
\begin{enumerate}
  \item Acyclicity $\Rightarrow$ $\mathrm{corr}(f) = -1/3$
        (Lemma~\lean{ArrowTheorem.acyclic_implies_corrFunc}).
  \item $\mathrm{corr}(f) = -1/3$ $\Rightarrow$ $\hat f(S) = 0$ for $|S|\ne 1$
        (Lemma~\lean{ArrowTheorem.corrFunc_eq_neg_third_of_weight_one}).
  \item Degree-1 + unanimous + $\pm 1$-valued $\Rightarrow$ dictator
        (Lemma~\lean{ArrowTheorem.degree_one_implies_dictator}).
\end{enumerate}
\end{theorem}
